\section{Temario}
Temario obtenido de la anterior convocatoria

\renewcommand{\thesubsubsection}{\thesubsection.\Roman{subsubsection}}
\subsection{A. Temas generales}
	\subsubsection{Marco constitucional español y Unión Europea}
	\begin{enumerate}
		\item La Constitución Española de 1978 (I). Los principios constitucionales. Los derechos fundamentales y sus garantías. La Corona. Cortes Generales. Congreso de los Diputados y Senado. El Gobierno. Los órganos constitucionales de control del Gobierno: Tribunal de Cuentas, Defensor del Pueblo. La función consultiva: el Consejo de Estado.
		\item La Constitución Española de 1978 (II). El Poder Judicial. La justicia en la Constitución. El Consejo General del Poder Judicial. El Ministerio Fiscal. El Tribunal Constitucional. Naturaleza, organización y atribuciones. 
		\item La Constitución Española de 1978 (III). Las Comunidades Autónomas: organización política y administrativa. La Administración local: regulación constitucional y entidades que la integran.
		\item La Constitución Española de 1978 (IV). La Administración pública: Principios constitucionales. La Administración General del Estado. Sus órganos centrales. 
		\item La Unión Europea: antecedentes, evolución y objetivos. Los tratados originarios y modificativos. El derecho de la Unión Europea. Relación entre el derecho de la Unión Europea y el ordenamiento jurídico de los Estados miembros. Las instituciones: el Consejo Europeo, el Consejo, la Comisión Europea, el Parlamento Europeo el Tribunal de Justicia de la Unión Europea, el Tribunal de Cuentas y el Banco Central Europeo.
		\item La organización territorial de la Administración General del Estado. Delegados y subdelegados del Gobierno. Organización de los servicios periféricos. El sector público institucional. Organización, funcionamiento y principios generales de actuación. Los organismos y entidades públicas estatales. 
	\end{enumerate}
	\subsubsection{Actuación administrativa y gestión financiera}
	\begin{enumerate}
		\setcounter{enumi}{6}
		\item Las fuentes del derecho administrativo. La jerarquía de las fuentes. La ley. Las disposiciones del Gobierno con fuerza de ley: decreto-ley y decreto legislativo. El reglamento: concepto, clases y límites. Otras fuentes del derecho administrativo.
		\item El régimen jurídico de las Administraciones públicas. El procedimiento administrativo y la relación de los ciudadanos con las Administraciones públicas.
		\item  La revisión de los actos en vía administrativa. Los recursos administrativos. El recurso contencioso-administrativo.
		\item La responsabilidad patrimonial de las Administraciones públicas. Procedimiento de responsabilidad patrimonial. La acción de responsabilidad.
		\item Los contratos de las Administraciones públicas. Principios comunes. La Ley de Contratos del Sector Público. Requisitos necesarios para la celebración de los contratos. Procedimientos de contratación y formas de adjudicación. Tipos de contratos y características generales. La facturación en el sector público.
		\item Los presupuestos generales del Estado. Estructura del presupuesto: clasificación de los gastos. Las modificaciones presupuestarias. La ejecución del gasto y su control interno: la Intervención General de la Administración del Estado. El control externo: el Tribunal de Cuentas y las Cortes Generales.
		\item Los convenios. Las encomiendas de gestión o encargos a medios propios personificados. Límites en su utilización. La colaboración público-privada en la prestación de los servicios públicos. El régimen jurídico de las subvenciones públicas. Procedimiento. 
	\end{enumerate}
	\subsubsection{Estructura social y económica de España}
	\begin{enumerate}
		\setcounter{enumi}{13}
		\item El modelo económico español en el marco de la economía mundial. El crecimiento sostenible y las políticas ambientales. La Agenda 2030 y los Objetivos de Desarrollo Sostenible. Líneas generales de la política económica actual. La renta nacional. Evolución y distribución. La internacionalización  de la economía española. Comercio exterior y balanza de pagos.
		\item El sector primario. El sector industrial. El sector servicios. Los subsectores de comercio, transportes y comunicaciones. El turismo.
		\item La estructura social. Tendencias demográficas. Políticas de igualdad de género. Discapacidad y dependencia. Políticas contra la violencia de género.
		\item El mercado de trabajo: evolución y características actuales. Principales magnitudes. Políticas públicas de empleo.
		\item El impacto de las tecnologías de la información y las comunicaciones en la sociedad. La brecha digital. Nuevos hábitos de relación y de consumo de la información.
		\item El impacto de las tecnologías de la información y las comunicaciones en la economía y en el mercado de trabajo. Principales magnitudes.
	\end{enumerate}
	\subsubsection{Dirección pública}
	\begin{enumerate}
		\setcounter{enumi}{19}
		\item La función gerencial en las Administraciones públicas. Particularidades de la gestión pública.
		\item Liderazgo. Gestión de competencias y personas. Gestión de conflictos e inteligencia emocional en el puesto de trabajo. Técnicas de negociación y gestión comercial.
		\item La calidad en los servicios públicos y el papel de las TIC en su modernización. La legislación en materia de sociedad de la información y administración electrónica en España y Europa.
		\item La crisis de la burocracia como sistema de gestión. La adecuación de técnicas del sector privado. El «management» público y la gobernanza. Hacia la excelencia en las instituciones públicas.
		\item Empleo Público en España. Tipología del empleo público. Estatuto Básico del Empleado Público. Planificación de recursos humanos en las Administraciones Públicas: la oferta de empleo público. Ley de Incompatibilidades del personal al servicio de las Administraciones Públicas.
		\item La gobernanza pública y el gobierno abierto. Concepto y principios informadores del gobierno abierto: colaboración, participación, transparencia y rendición de cuentas. Datos abiertos y reutilización. El marco jurídico y los planes de gobierno abierto en España. La Ley 19/2013, de 9 de diciembre, de transparencia, acceso a la información pública y buen gobierno. El Consejo de Transparencia y Buen Gobierno: Real Decreto 919/2014, de 31 de octubre, por el que se aprueba su estatuto. Funciones. La Oficina de Transparencia y Acceso a la Información (OTAI): Funciones. El Portal de Transparencia. Las Unidades de Información y Transparencia (UITS): Funciones. La transparencia y el acceso a la información en las comunidades autónomas y entidades locales.
		\item La dinamización y el apoyo a la actividad económica y al emprendimiento. La garantía de la unidad de mercado.
		\item La política de protección de datos de carácter personal. Régimen jurídico. El Reglamento UE 2016/679, de 27 de abril, relativo a la protección de las personas físicas en lo que respecta al tratamiento de datos personales y a la libre circulación de estos datos. Principios y derechos. Obligaciones. El Delegado de Protección de Datos en las Administraciones Públicas. La Agencia Española de Protección de Datos.
	\end{enumerate}
\subsection{B. Temas específicos}
	\subsubsection{Organización y gestión de los sistemas de información}
	\begin{enumerate}
		\setcounter{enumi}{27}
		\item Definición, estructura, y dimensionamiento eficiente de los sistemas de información.
		\item La información en las organizaciones. Las organizaciones basadas en la información. La Administración como caso específico de este tipo de organización.
		\item Modelos de gobernanza TIC. Organización e instrumentos operativos de las tecnologías de la información y las comunicaciones en la Administración General del Estado y sus organismos públicos. La transformación digital de la Administración General del Estado.
		\item Reutilización de la información en el sector público en Europa y España. Papel de las TIC en la implantación de políticas de datos abiertos y transparencia.
		\item Estrategia, objetivos y funciones del directivo de sistemas y tecnologías de la información en la Administración.
		\item Herramientas de planificación y control de gestión de la función del directivo de sistemas y tecnologías de la información en la Administración. El cuadro de mando.
		\item Organización y funcionamiento de un centro de sistemas de información. Funciones de desarrollo, mantenimiento, sistemas, bases de datos, comunicaciones, seguridad, calidad, microinformática y atención a usuarios.
		\item Dirección y gestión de proyectos de tecnologías de la información. Planificación estratégica, gestión de recursos, seguimiento de proyectos, toma de decisiones.
		\item Metodologías predictivas para la gestión de proyectos: GANTT, PERT.
		\item Metodologías ágiles para la gestión de proyectos. Metodologías lean.
		\item Auditoría informática. Concepto y contenidos. Administración, planeamiento, organización, infraestructura técnica y prácticas operativas.
		\item La gestión de la compra pública de tecnologías de la información.
		\item Adquisición de sistemas: estudio de alternativas, evaluación de la viabilidad y toma de decisión.
		\item Alternativas básicas de decisión en el campo del equipamiento hardware y software.
		\item La rentabilidad de las inversiones en los proyectos de tecnologías de la información.
		\item La protección jurídica de los programas de ordenador. Los medios de comprobación de la legalidad y control del software.
		\item Accesibilidad y usabilidad. W3C. Diseño universal. Diseño web adaptativo. 
		\item Interoperabilidad de sistemas (1). El Esquema Nacional de Interoperabilidad. Dimensiones de la interoperabilidad.
		\item Interoperabilidad de sistemas (2). Las Normas Técnicas de Interoperabilidad. Interoperabilidad de los documentos y expedientes electrónicos y normas para el intercambio de datos entre Administraciones Públicas.
		\item Seguridad de sistemas (1). Análisis y gestión de riesgos. Herramientas.
		\item Seguridad de sistemas (2). El Esquema Nacional de Seguridad. Adecuación al Esquema Nacional de Seguridad. Estrategia Nacional de Seguridad. CCN-STIC.
		\item Infraestructuras, servicios comunes y compartidos para la interoperabilidad entre Administraciones públicas. Cl@ve, la carpeta ciudadana, el Sistema de Interconexión de Registros, la Plataforma de Intermediación de Datos, y otros servicios.
		\item Organizaciones internacionales y nacionales de normalización. Pruebas de conformidad y certificación. El establecimiento de servicios de pruebas de conformidad.
		\item Planes y Actuaciones de la Agenda Digital para España. Descripción, estructura y objetivos de los planes. El Mercado Único Digital.
	\end{enumerate}
	\subsubsection{Tecnología básica}
	\begin{enumerate}
		\setcounter{enumi}{51}
		\item Sistemas de altas prestaciones. Grid Computing. Mainframe.
		\item Equipos departamentales. Servidores. Medidas de seguridad para equipos departamentales y servidores. Centros de proceso de datos: diseño, implantación y gestión.
		\item Dispositivos personales de PC y dispositivos móviles. La conectividad de los dispositivos personales. Medidas de seguridad y gestión para equipos personales y dispositivos móviles.
		\item Cloud Computing. IaaS, PaaS, SaaS. Nubes privadas, públicas e híbridas.
		\item Sistemas de almacenamiento para sistemas grandes y departamentales. Dispositivos para tratamiento de información multimedia. Virtualización del almacenamiento. Copias de seguridad.
		\item Tipos de sistemas de información multiusuario. Sistemas grandes, medios y pequeños. Servidores de datos y de aplicaciones. Virtualización de servidores.
		\item El procesamiento cooperativo y la arquitectura cliente-servidor. Arquitectura SOA.
		\item Conceptos y fundamentos de sistemas operativos. Evolución y tendencias.
		\item Sistemas operativos UNIX-LINUX. Fundamentos, administración, instalación, gestión.
		\item Sistemas operativos Microsoft. Fundamentos, administración, instalación, gestión.
		\item Conceptos básicos de otros sistemas operativos: OS X, iOS, Android, z/OS. Sistemas operativos para dispositivos móviles.
		\item Los sistemas de gestión de bases de datos SGBD. El modelo de referencia de ANSI.
		\item El modelo relacional. El lenguaje SQL. Normas y estándares para la interoperabilidad entre gestores de bases de datos relacionales.
		\item Arquitectura de desarrollo en la web. Desarrollo web front-end. Scripts de cliente. Frameworks. UX. Desarrollo web en servidor, conexión a bases de datos e interconexión con sistemas y servicios.
		\item Entorno de desarrollo Microsoft.NET.
		\item Entorno de desarrollo JAVA.
		\item Entorno de desarrollo PHP.
		\item Software de código abierto. Software libre. Conceptos base. Aplicaciones en entorno ofimático y servidores web.
		\item Inteligencia artificial: la orientación heurística, inteligencia artificial distribuida, agentes inteligentes.
		\item Gestión del conocimiento: representación del conocimiento. Sistemas expertos. Herramientas.
		\item Sistemas CRM (Customer Relationship Management) y ERP (Enterprise Resource Planning). Generación de informes a la dirección.
		\item E-learning: conceptos, herramientas, sistemas de implantación y normalización. 
		\item Los sistemas de información geográfica. Conceptos y funcionalidad básicos.
		\item Gestión de los datos corporativos. Almacén de datos (Data-Warehouse). Arquitectura OLAP. Minería de datos.
		\item Big Data. Captura, análisis, transformación, almacenamiento y explotación de conjuntos masivos de datos. Entornos Hadoop o similares. Bases de datos NoSQL.
		\item Lenguajes y herramientas para la utilización de redes globales. HTML, CSS y XML. Navegadores web y compatibilidad con estándares.
		\item Comercio electrónico. Mecanismos de pago. Gestión del negocio. Factura electrónica. Pasarelas de pago.
		\item El cifrado. Algoritmos de cifrado simétricos y asimétricos. La función hash. El notariado.
		\item Identificación y firma electrónica (1) Marco europeo y nacional. Certificados digitales. Claves privadas, públicas y concertadas. Formatos de firma electrónica. Protocolos de directorio basados en LDAP y X.500. Otros servicios.
		\item Identificación y firma electrónica (2) Prestación de servicios públicos y privados. Infraestructura de clave pública (PKI). Mecanismos de identificación y firma: «Smart Cards», DNI electrónico, mecanismos biométricos.
		\item Adaptación de aplicaciones y entornos a los requisitos de la normativa de protección de datos según los niveles de seguridad. Herramientas de cifrado y auditoría.
		\item El tratamiento de imágenes. Tecnologías de digitalización y de impresión. Impresión 3D.
		\item Reconocimiento óptico de caracteres (OCR, ICR). Reconocimiento biométrico.
	\end{enumerate}
	\subsubsection{Ingeniería de los sistemas de información}
	\begin{enumerate}
		\setcounter{enumi}{84}
		\item El ciclo de vida de los sistemas de información. Modelos del ciclo de vida. 86. Planificación estratégica de sistemas de información y de comunicaciones. El plan de sistemas de información.
		\item Análisis funcional de sistemas, casos de uso e historias de usuario. Metodologías de desarrollo de sistemas. Metodologías ágiles: Scrum y Kanban.
		\item Análisis del dominio de los sistemas: modelado de dominio, modelo entidad relación y modelos de clases.
		\item Análisis dinámico de sistemas: modelado de procesos, modelado dinámico y BPMN (Business Process Model and Notation).
		\item Análisis de aspectos no funcionales: rendimiento, seguridad, privacidad.
		\item Diseño arquitectónico de sistemas. Diagramas de despliegue.
		\item Técnicas de diseño de software. Diseño por capas y patrones de diseño.
		\item La elaboración de prototipos en el desarrollo de sistemas. Diseño de interfaces de aplicaciones.
		\item La metodología de planificación y desarrollo de sistemas de información Métrica.
		\item Procesos de pruebas y garantía de calidad en el desarrollo de software. Planificación, estrategia de pruebas y estándares. Niveles, técnicas y herramientas de pruebas de software. Criterios de aceptación de software.
		\item Modelos de integración continua. Herramientas y sus aplicaciones.
		\item Métricas y evaluación de la calidad del software. La implantación de la función de calidad.
		\item La estimación de recursos y esfuerzo en el desarrollo de sistemas de información.
		\item La migración de aplicaciones en el marco de procesos de ajuste dimensional y por obsolescencia técnica. Gestión de la configuración y de versiones. Gestión de entornos.
		\item Mantenimiento de sistemas. Mantenimiento predictivo, adaptativo y correctivo. Planificación y gestión del mantenimiento.
		\item Gestión de cambios en proyectos de desarrollo de software. Gestión de la configuración y de versiones. Gestión de entornos.
		\item La calidad en los servicios de información. El Modelo EFQM y la Guía para los servicios ISO 9004.
		\item Gestión documental. Gestión de contenidos. Tecnologías CMS y DMS de alta implantación.
		\item Sistemas de recuperación de la información. Políticas, procedimientos y métodos para la conservación de la información.
		\item Planificación y control de las TIC: gestión de servicios e infraestructuras TIC, gestión del valor de las TIC. Acuerdos de nivel de servicio. Gestión de incidencias. Bases conceptuales de ITIL (IT Infrastructure Library), y CoBIT (Control Objetives for Information and Related Technology), objetivos de control y métricas.
		\item Las tecnologías emergentes. Concepto. Clasificación, aspectos jurídicos y aplicaciones.
	\end{enumerate}
	\subsubsection{Redes, comunicaciones e Internet}
	\begin{enumerate}
		\setcounter{enumi}{105}
		\item Redes de telecomunicaciones. Conceptos. Medios de transmisión. Conmutación de circuitos y paquetes. Protocolos de encaminamiento. Infraestructuras de acceso. Interconexión de redes. Calidad de servicio.
		\item La red Internet y los servicios básicos.
		\item Sistemas de cableado y equipos de interconexión de redes.
		\item El modelo de referencia de interconexión de sistemas abiertos (OSI) de ISO: arquitectura, capas, interfaces, protocolos, direccionamiento y encaminamiento.
		\item Tecnologías de acceso: fibra (GPON, FTTH), móviles (LTE), inalámbrica. 
		\item Redes de transporte: JDSxWDM, MPLS. Redes de agregación: ATM, CarrierEthernet-VPLS (H-VPLS).
		\item Redes inalámbricas: el estándar IEEE 802.11.Características funcionales y técnicas. Sistemas de expansión del espectro. Sistemas de acceso. Autenticación. Modos de operación. Bluetooth. Seguridad, normativa reguladora.
		\item Redes IP: arquitectura de redes, encaminamiento y calidad de servicio. Transición y convivencia IPv4 – IPv6. Funcionalidades específicas de IPv6.
		\item Redes de nueva generación y servicios convergentes (NGN/IMS). VoIP, ToIP y comunicaciones unificadas.
		\item La transformación digital e industria 4.0: ciudades inteligentes. Internet de las Cosas (IoT).
		\item Redes de área local. Arquitectura. Tipología. Medios de transmisión. Métodos de acceso. Dispositivos de interconexión. Gestión de dispositivos. Administración de redes LAN. Gestión de usuarios en redes locales. Monitorización y control de tráfico. Gestión SNMP. Configuración y gestión de redes virtuales (VLAN). Redes de área extensa. 
		\item Arquitectura de las redes Intranet y Extranet. Concepto, estructura y características. Su implantación en las organizaciones. Modelo de capas: servidores de aplicaciones, servidores de datos, granjas de servidores. 
		\item Las redes públicas de transmisión de datos. La red SARA. La red sTESTA. Planificación y gestión de redes.
		\item Telecomunicaciones por cable (CATV). Estructura de la red de cable. Operadores del mercado. Servicios de red.
		\item El correo electrónico. Servicios de mensajería. Servicios de directorio.
		\item Las comunicaciones móviles. Generaciones de tecnologías de telefonía móvil.
		\item Aplicaciones móviles. Características, tecnologías, distribución y tendencias.
		\item La seguridad en redes. Tipos de ataques y herramientas para su prevención: cortafuegos, control de accesos e intrusiones, técnicas criptográficas, etc. Medidas específicas para las comunicaciones móviles.
		\item La seguridad en el nivel de aplicación. Tipos de ataques y protección de servicios web, bases de datos e interfaces de usuario.
		\item Ciberseguridad. La estrategia nacional de ciberseguridad.
		\item La gestión de la continuidad del negocio. Planes de Continuidad y Contingencia del Negocio.
		\item Normas reguladoras de las telecomunicaciones. La Comisión Nacional de los Mercados y la Competencia  CNMC): organización, funciones y competencia en el ámbito de las telecomunicaciones.
		\item Sistemas de videoconferencia. Herramientas de trabajo en grupo. Dimensionamiento y calidad de servicio en las comunicaciones y acondicionamiento de salas y equipos. Streaming de video.
		\item Acceso remoto a sistemas corporativos: gestión de identidades, single sign-on y teletrabajo.
		\item Virtualización de sistemas y de centros de datos. Virtualización de puestos de trabajo. Maquetas de terminales Windows y de servidores Linux.
		\item Herramientas de trabajo colaborativo y redes sociales. La Guía de comunicación digital de la Administración del Estado.
	\end{enumerate}


\renewcommand{\thesubsubsection}{\thesubsection.\arabic{subsubsection}}